\section{Deployment Guidance for SOC Teams}

\ARES{} translates into concrete guidance for teams deploying
autonomous AI agents, organised around six questions SOC architects
face when integrating agentic AI, and closing with a five-step
pre-deployment checklist.

\textbf{When should AI agents be trusted?}
ARS is proposed as a pre-execution screening signal: an agent scoring
below a learned threshold is escalated to a human or reconfigured
(add knowledge context, restore tool access). The threshold is
deployment-specific and must be fit on the deployer's own labelled
outcomes, not transplanted (our analytical corpus is illustrative
only). Because four of the five dimensions ($\delta_S, \delta_K,
\delta_C, \delta_T$) are computable before execution, screening adds no
runtime overhead and slots into SOAR platforms as a pre-condition for
autonomous execution~\cite{nist_ai_rmf}. ARS is a screening heuristic,
not a validated detector; the calibration signal below is the
empirically supported lever.

\textbf{How can reliability be monitored continuously?}
The Epistemic Gap enables runtime hallucination risk monitoring: by
comparing the agent's stated confidence against its known knowledge
coverage, a monitor flags high-risk outputs for human review before
they trigger automated actions, valuable for host isolation, threat
feed updates, and evidence preservation. Computing $\mathcal{E}$
requires only the agent's confidence and pre-computed $\delta_K$,
both already in the execution record.

\textbf{How should knowledge maintenance be scheduled?}
The temporal-drift model gives the \textit{shape} of a per-domain
update schedule: reliability decays at a domain-specific rate
$\lambda$, so fast-moving domains need more frequent refresh.
Illustrative constants ($\lambda = 0.005$~day$^{-1}$ for ransomware,
$0.001$ for APT attribution) imply roughly quarterly versus annual
refresh, but must be fit on each deployer's own drift data; the
contribution is the schedule form, not these numbers.

\textbf{How should multi-agent systems be designed?}
Composition never beat the best single agent live, but the penalty was
small ($3.3$~pp) and model-dependent, not the large amplification the
model predicts. Treat decomposition as a reliability cost to justify,
not a default: measure the depth penalty for your own backbone. Where
composition is warranted, three mitigations reduce cascade risk: (1)
inter-agent verification checkpoints; (2) confidence-gated propagation
(low-confidence outputs go to humans, not downstream); (3) redundant
agent paths for critical decisions.

\textbf{How can safer agent systems be built?}
Our calibration-gap result reorders the usual priorities. (1)
\textit{Calibration-first design}, the operative lever: prefer and
validate agents whose expressed confidence falls as their knowledge
degrades, since it is calibration, not knowledge coverage alone, that
determines whether a knowledge gap becomes a confident fabrication.
(2) \textit{Knowledge coverage} still matters, but chiefly for
\textit{task success}: a well-calibrated but under-informed agent will
abstain safely, trading coverage for reliability, so knowledge refresh
is a capability investment more than a safety one. (3) \textit{Fail-safe
defaults}, preserve and reward the abstention behaviour real agents
already exhibit, rather than prompting them to always answer. The one
agent configured or fine-tuned to answer confidently under uncertainty
is the dangerous exception, and the one to screen out.

\textbf{How does this map to regulatory and governance frameworks?}
The NIST AI Risk Management
Framework~\cite{nist_ai_rmf} identifies ``valid and reliable,''
``accountable and transparent,'' and ``safe'' as top-tier
trustworthiness characteristics: \ARES{} operationalises the first
via ARS, the second via the traceable mismatch decomposition, and
the third via the Epistemic-Gap escalation gate. At the application layer,
the OWASP Top~10 for Agentic
Applications~\cite{owasp2026agentic} catalogues capability
over-provisioning, tool misuse, and confident hallucination as
distinct risk categories, each of which has a direct analogue in
\ARES{}: $\delta_C$ over-provisioning, $\delta_C$ plus $\delta_R$
tool misuse, and the Epistemic Gap for confident hallucination.

\textbf{A deployment checklist.}
Distilling the above into a pragmatic checklist for SOC
architects yields five steps: \textit{(i)} estimate pre-execution
$\delta_S$, $\delta_K$, $\delta_C$, $\delta_T$ for each intended
task class; \textit{(ii)} apply the task-class-specific weight
vector and compare ARS against a threshold fit on local labelled
outcomes; \textit{(iii)}
if ARS is below threshold, choose one of three mitigations
(knowledge refresh, capability provisioning, or task-scope
restriction) and re-score; \textit{(iv)} instrument the runtime
to emit stated confidence alongside each decision and compute
$\mathcal{E}$ per record, and prefer agents whose confidence
falls under withheld knowledge; \textit{(v)} schedule knowledge
refreshes according to a domain-specific temporal-drift
constant fit on local drift data. Each step maps to a concrete
engineering artefact, so the reliability gate does not stay
theoretical.
